\documentclass[10pt,a4paper]{amsart}
\usepackage[margin=19mm]{geometry}
\usepackage{amsmath,amssymb,mathtools}
\usepackage[scr=rsfs,cal=euler]{mathalfa}
\usepackage{xfrac}
\usepackage[unicode,hidelinks,bookmarks=false]{hyperref}
\newcommand{\ie}{i.e.\ }
\newcommand{\cf}{cf.\ }
\newcommand{\resp}{respectively\ }
\newcommand{\Subsection}{\subsection}
\providecommand{\nobreakdash}{\nobreak\hbox{-}}
\setlength{\emergencystretch}{2em}
\multlinegap0pt
\usepackage{amssymb}
\usepackage{mathtools}
\newtheorem{definition}{Definition}
\newtheorem{proposition}{Proposition}
\title{An Algebraic Proof of the Goldbach Conjecture}
\author{Jason R. South}
\date{Source: arXiv:2206.01179v16 (CC BY 4.0)}
\begin{document}
\maketitle

\noindent\emph{Source and licence.} An Algebraic Proof of the Goldbach Conjecture. Version arXiv:2206.01179v16 (CC BY 4.0), distributed by arXiv under the Creative Commons Attribution 4.0 International licence, http://creativecommons.org/licenses/by/4.0/, as recorded in the arXiv licence metadata for this exact version and shown on the abstract page https://arxiv.org/abs/2206.01179v16. Full licence notice: \texttt{licenses/arxiv-2206.01179-CC-BY-4.0.txt}.\par\smallskip

\noindent\textit{Editorial scope.} Counted unit: the article's proposition gbf (source0.tex lines 201--203). The article's complete original proof of it is delivered with it (source0.tex lines 205--207, one \texttt{proof} environment of 3 lines). No mathematics is written, repaired or re-derived: the statement and the proof are the article's own lines, in the article's own notation, and no retained line is substituted or re-flowed.  Retained uncounted as the source-local context the proof needs: lines 164--177; lines 174--176.  Nothing was omitted from any retained span, so the package carries no omission record.  No retained line contains a citation, so the wrapper carries no bibliography and no bibliography entry.  No retained line contains a figure environment, an image-inclusion command or drawing code, so no image is delivered and the package declares no asset dependency.  Editorial formatting only, in addition: an AMS \texttt{amsart} preamble that restates the article's own declarations and macros used by the retained lines, namely the environments \texttt{definition}, \texttt{proposition} on separate counters, together with the standard packages those lines need.  The retained environments print in this document as Definition 1, Proposition 1 in order of appearance, whereas the article prints the same items with the numbering of its own full document.  The complete native source of the article as distributed in the arXiv e-print for this exact version is bundled unchanged as \texttt{sources/126459/source0.tex}.
\begin{definition}\label{polyx}
A \textit{Goldbach Polynomial} (G.P.) exists when it is possible to take some $a \in \mathbb{N}$ and for each $p_i \in \mathcal{P}_a$ a unique $\alpha_i \in \mathbb{N} \cup \{0\}$ exists producing a mapping $\mathcal{G}_- : \mathbb{C} \to \mathbb{C}$ where  
\begin{equation}\label{roots}
 \mathcal{G}_-(z) = \prod_{p_k \in \mathcal{P}_a} (z - p_k) - \prod_{p_k \in \mathcal{P}_a}p_k^{\alpha_k}
 \end{equation}
 with the condition $\mathcal{G}_-(2a) = 0$. Let the set of roots for the G.P. be given by 
\begin{equation}\label{R}
\mathcal{R}_a = \{ r_i \in \mathbb{C} : \mathcal{G}_-(r_i) = 0 \}
\end{equation}
the factorizations below will be needed in future sections. 
 \begin{equation}\label{factor1}
 \prod_{p_k \in \mathcal{P}_a}(r_i - p_k) = \prod_{p_k \in \mathcal{P}_a}p_k^{\alpha_k} \quad \text{where} \quad -(p_i - r_i) \prod_{p_k \in \mathcal{P}_a \setminus\{p_i\}}(r_i - p_k) = -\prod_{p_k \in \mathcal{P}_a}p_k^{\alpha_k}.
 \end{equation}
\end{definition}

 \begin{equation}
 \prod_{p_k \in \mathcal{P}_a}(r_i - p_k) = \prod_{p_k \in \mathcal{P}_a}p_k^{\alpha_k} \quad \text{where} \quad -(p_i - r_i) \prod_{p_k \in \mathcal{P}_a \setminus\{p_i\}}(r_i - p_k) = -\prod_{p_k \in \mathcal{P}_a}p_k^{\alpha_k}.
 \end{equation}

\begin{proposition}\label{gbf}
\textit{Let $a \in \mathbb{N}_{> 3}$. If $2a$ is a counterexample to the G.C. then there exists a G.P.}
\end{proposition}

\begin{proof}
Assume $2a$ is counter-example to the G.C. Then no $2a - p_k$ may be prime when $p_k \in \mathcal{P}_a$. If any $2a - p_k$ had a prime factor $a < p < 2a$ then $0< 2a - p_k = n p < 2a$ with $p > a$ forcing $0<n < 2$. Hence, $2a - p_k = p$, contradicting the assumption that $2a$ is a counter-example to the G.C. Therefore, any counter-example requires each $2a - p_k$ only to be divisible by primes in $\mathcal{P}_a$ where the Fundamental Theorem of Arithmetic ensures unique $\alpha_1, \dots, \alpha_{\pi(a)} \in \mathbb{N}\cup \{0\}$ exist satisfying equation \ref{factor1}. The existence of these $\alpha_i$ produce a G.P. from Definition \ref{polyx}.
\end{proof}

\end{document}
